\documentclass[11pt,a4paper]{article}
% main (standalone preamble for editing)
% vim: nowrap: spell:
% ══════════════════════════════════════════════════════════════════════════════
% Speed Maths --- shared preamble
% Included by every sheet and answer file via:  % main (standalone preamble for editing)
% vim: nowrap: spell:
% ══════════════════════════════════════════════════════════════════════════════
% Speed Maths --- shared preamble
% Included by every sheet and answer file via:  % main (standalone preamble for editing)
% vim: nowrap: spell:
% ══════════════════════════════════════════════════════════════════════════════
% Speed Maths --- shared preamble
% Included by every sheet and answer file via:  \input{../../shared/preamble}
% Editing THIS file changes the look of every sheet/answer across every pillar.
% ══════════════════════════════════════════════════════════════════════════════

\usepackage[a4paper, top=25mm, bottom=25mm, left=22mm, right=22mm]{geometry}
\usepackage{amsmath, amssymb}
\usepackage{enumitem}
\usepackage{booktabs}
\usepackage{array}
\usepackage{parskip}
\usepackage{microtype}
\usepackage{fancyhdr}
\usepackage{titlesec}
\usepackage{mdframed}
\usepackage{xcolor}
\usepackage[hidelinks]{hyperref}
\usepackage{graphicx}
\usepackage{tikz}
\usepackage{pgfplots}
\pgfplotsset{compat=1.18}

% ── Colours ───────────────────────────────────────────────────────────────────
\definecolor{darkgray}{gray}{0.15}
\definecolor{medgray}{gray}{0.45}
\definecolor{lightgray}{gray}{0.93}
\definecolor{boxgray}{gray}{0.88}

% ── Section formatting ──────────────────────────────────────────────────────────
\titleformat{\section}[block]
  {\normalfont\large\bfseries}{}
  {0pt}{}[\vspace{2pt}\hrule\vspace{6pt}]
\titlespacing*{\section}{0pt}{14pt}{4pt}

% ── List formatting ─────────────────────────────────────────────────────────────
\setlist[enumerate,1]{label=\textbf{\arabic*.}, leftmargin=2em, itemsep=10pt}

% ── Branding constants (edit here, not per-file) ────────────────────────────────
\newcommand{\SpeedExamLine}{\textbf{Speed Maths:} Competition \& Exam Prep}
\newcommand{\SpeedCredit}{Series by \href{https://www.linkedin.com/in/cerealdev/}{David Akinyele-Aje}}

% ── Header / footer ──────────────────────────────────────────────────────────────
% Usage (once, after \input{../../shared/preamble} and before \begin{document}):
%   \SpeedHeader{Algebra}{3}
% First arg  = pillar name shown as "Speed <Pillar> --- Daily Drill"
% Second arg = worksheet number
\pagestyle{fancy}
\newcommand{\SpeedHeader}[2]{%
  \fancyhf{}%
  \renewcommand{\headrulewidth}{0.4pt}%
  \setlength{\headheight}{14pt}%
  \fancyhead[L]{\small\textbf{Speed #1 --- Daily Drill}}%
  \fancyhead[R]{\small Worksheet \##2}%
  \fancyfoot[L]{\small \SpeedExamLine}%
  \fancyfoot[C]{\small\thepage}%
  \fancyfoot[R]{\small \SpeedCredit}%
  \renewcommand{\footrulewidth}{0.4pt}%
}

% ── Title block (used at the top of every sheet AND answer file) ───────────────
% Usage: \SpeedTitleBlock{Daily Algebra Drill \#3}{Jane Doe}
% %% AGENTS: When generating a new sheet, ask the human contributor for the name they want credited in argument 2.
% For answer files, pass the "--- Answers \& Investigations" suffix in the arg.
\newcommand{\SpeedTitleBlock}[2]{%
  \begin{center}
    {\LARGE\bfseries #1}\\[4pt]
    {\normalsize\itshape Sheet by #2}\\[6pt]
    \hrule\vspace{4pt}
    \begin{tabular}{llcc}
      \toprule
      \textbf{Section} & \textbf{Topic} & \textbf{Questions} & \textbf{Target time}\\
      \midrule
      A & Rapid Recognition          & 10 & 2 min 30 sec \\
      B & Manipulation Drills        & 10 & 8 min        \\
      C & Substitution \& Structure  & 8  & 10 min       \\
      D & Challenge                  & 5  & 15 min       \\
      \bottomrule
    \end{tabular}
    \vspace{4pt}
    \hrule
  \end{center}
}

% ── Sheet metadata: topic + the day's new tools ────────────────────────────────
% Usage (immediately after \SpeedTitleBlock, sheets only):
%   \SpeedMeta{Expanding, factorising, and the identity toolkit}
%             {difference of squares, sum/product form, completing the square}
%
% Arg 1 = the sheet's topic in one line. The website lists this instead of
%         "Sheet 03", so write the thing a reader would actually search for.
% Arg 2 = comma-separated, at most five: the tools this sheet introduces that
%         earlier sheets in the pillar did not. These become the website's tags.
%         Leave out anything the reader already met — the tags are meant to say
%         what is new today, not to summarise the sheet.
%
% Both are parsed straight out of this file by tools/build_website.py, so the
% sheet stays the single source of truth and the site cannot drift from it.
%
% This renders NOTHING. It is a declaration for the build script, not a piece of
% the sheet's layout: adding it to a sheet must not change that sheet's PDF, so
% the 25 existing sheets can be annotated without a single one of them being
% reissued. If the toolkit should also appear on the page, that is a separate
% decision about the sheet design — make it deliberately, here, not by leaving
% this macro printing something nobody asked it to.
\newcommand{\SpeedMeta}[2]{}

% ── Answer / Method / Investigate-further boxes (answer files only) ────────────
\newmdenv[
  backgroundcolor=lightgray,
  linecolor=medgray,
  linewidth=0.5pt,
  leftmargin=0pt, rightmargin=0pt,
  innerleftmargin=8pt, innerrightmargin=8pt,
  innertopmargin=5pt, innerbottommargin=5pt,
  skipabove=4pt, skipbelow=4pt
]{ansbox}

\newmdenv[
  backgroundcolor=white,
  linecolor=darkgray,
  linewidth=0.8pt,
  leftline=true, rightline=false, topline=false, bottomline=false,
  leftmargin=0pt, rightmargin=0pt,
  innerleftmargin=8pt, innerrightmargin=6pt,
  innertopmargin=4pt, innerbottommargin=4pt,
  skipabove=4pt, skipbelow=6pt
]{invbox}

\newcommand{\ans}[1]{%
  \begin{ansbox}
    \textbf{Answer:} #1
  \end{ansbox}}

\newcommand{\method}[1]{%
  {\small\textbf{Method:} #1}\par}

\newcommand{\inv}[1]{%
  \begin{invbox}
    \textbf{\small Investigate further:} {\small #1}
  \end{invbox}}

% ── Misc helpers ────────────────────────────────────────────────────────────────
\newcommand{\qn}[1]{\textbf{#1.}\;}

% Mistake-linking: attach a Top 5 Patterns / Common Traps bullet to the question
% label(s) in this sheet where it actually shows up, e.g. \seealso{B6, D2}
\newcommand{\seealso}[1]{\hfill\textit{\footnotesize(see #1)}}

% ── Graph options macro ─────────────────────────────────────────────────────────
% Usage: \GraphOptions{tikz code for A}{tikz code for B}{tikz code for C}{tikz code for D}
\newcommand{\GraphOptions}[4]{%
  \begin{center}
    \begin{tabular}{cc}
      \textbf{A)} & \textbf{B)} \\
      #1 & #2 \\[10pt]
      \textbf{C)} & \textbf{D)} \\
      #3 & #4
    \end{tabular}
  \end{center}
}

% ── Closing motivational rule + cross-promo footer (sheets only) ────────────────
% Usage: \SpeedClosing{"Quote text."}
\newcommand{\SpeedClosing}[1]{%
  \vspace{20pt}
  \begin{center}
  \rule{0.6\textwidth}{0.4pt}\\[4pt]
  {\small\itshape #1}\\[4pt]
  \rule{0.6\textwidth}{0.4pt}
  \end{center}
  \begin{center}
  {\small Found a mistake or want to contribute a sheet?}\\
  {\small Visit the open-source project at \href{https://github.com/The-CerealDev/speed-maths}{github.com/The-CerealDev/speed-maths}}
  \end{center}
}

% Editing THIS file changes the look of every sheet/answer across every pillar.
% ══════════════════════════════════════════════════════════════════════════════

\usepackage[a4paper, top=25mm, bottom=25mm, left=22mm, right=22mm]{geometry}
\usepackage{amsmath, amssymb}
\usepackage{enumitem}
\usepackage{booktabs}
\usepackage{array}
\usepackage{parskip}
\usepackage{microtype}
\usepackage{fancyhdr}
\usepackage{titlesec}
\usepackage{mdframed}
\usepackage{xcolor}
\usepackage[hidelinks]{hyperref}
\usepackage{graphicx}
\usepackage{tikz}
\usepackage{pgfplots}
\pgfplotsset{compat=1.18}

% ── Colours ───────────────────────────────────────────────────────────────────
\definecolor{darkgray}{gray}{0.15}
\definecolor{medgray}{gray}{0.45}
\definecolor{lightgray}{gray}{0.93}
\definecolor{boxgray}{gray}{0.88}

% ── Section formatting ──────────────────────────────────────────────────────────
\titleformat{\section}[block]
  {\normalfont\large\bfseries}{}
  {0pt}{}[\vspace{2pt}\hrule\vspace{6pt}]
\titlespacing*{\section}{0pt}{14pt}{4pt}

% ── List formatting ─────────────────────────────────────────────────────────────
\setlist[enumerate,1]{label=\textbf{\arabic*.}, leftmargin=2em, itemsep=10pt}

% ── Branding constants (edit here, not per-file) ────────────────────────────────
\newcommand{\SpeedExamLine}{\textbf{Speed Maths:} Competition \& Exam Prep}
\newcommand{\SpeedCredit}{Series by \href{https://www.linkedin.com/in/cerealdev/}{David Akinyele-Aje}}

% ── Header / footer ──────────────────────────────────────────────────────────────
% Usage (once, after % main (standalone preamble for editing)
% vim: nowrap: spell:
% ══════════════════════════════════════════════════════════════════════════════
% Speed Maths --- shared preamble
% Included by every sheet and answer file via:  \input{../../shared/preamble}
% Editing THIS file changes the look of every sheet/answer across every pillar.
% ══════════════════════════════════════════════════════════════════════════════

\usepackage[a4paper, top=25mm, bottom=25mm, left=22mm, right=22mm]{geometry}
\usepackage{amsmath, amssymb}
\usepackage{enumitem}
\usepackage{booktabs}
\usepackage{array}
\usepackage{parskip}
\usepackage{microtype}
\usepackage{fancyhdr}
\usepackage{titlesec}
\usepackage{mdframed}
\usepackage{xcolor}
\usepackage[hidelinks]{hyperref}
\usepackage{graphicx}
\usepackage{tikz}
\usepackage{pgfplots}
\pgfplotsset{compat=1.18}

% ── Colours ───────────────────────────────────────────────────────────────────
\definecolor{darkgray}{gray}{0.15}
\definecolor{medgray}{gray}{0.45}
\definecolor{lightgray}{gray}{0.93}
\definecolor{boxgray}{gray}{0.88}

% ── Section formatting ──────────────────────────────────────────────────────────
\titleformat{\section}[block]
  {\normalfont\large\bfseries}{}
  {0pt}{}[\vspace{2pt}\hrule\vspace{6pt}]
\titlespacing*{\section}{0pt}{14pt}{4pt}

% ── List formatting ─────────────────────────────────────────────────────────────
\setlist[enumerate,1]{label=\textbf{\arabic*.}, leftmargin=2em, itemsep=10pt}

% ── Branding constants (edit here, not per-file) ────────────────────────────────
\newcommand{\SpeedExamLine}{\textbf{Speed Maths:} Competition \& Exam Prep}
\newcommand{\SpeedCredit}{Series by \href{https://www.linkedin.com/in/cerealdev/}{David Akinyele-Aje}}

% ── Header / footer ──────────────────────────────────────────────────────────────
% Usage (once, after \input{../../shared/preamble} and before \begin{document}):
%   \SpeedHeader{Algebra}{3}
% First arg  = pillar name shown as "Speed <Pillar> --- Daily Drill"
% Second arg = worksheet number
\pagestyle{fancy}
\newcommand{\SpeedHeader}[2]{%
  \fancyhf{}%
  \renewcommand{\headrulewidth}{0.4pt}%
  \setlength{\headheight}{14pt}%
  \fancyhead[L]{\small\textbf{Speed #1 --- Daily Drill}}%
  \fancyhead[R]{\small Worksheet \##2}%
  \fancyfoot[L]{\small \SpeedExamLine}%
  \fancyfoot[C]{\small\thepage}%
  \fancyfoot[R]{\small \SpeedCredit}%
  \renewcommand{\footrulewidth}{0.4pt}%
}

% ── Title block (used at the top of every sheet AND answer file) ───────────────
% Usage: \SpeedTitleBlock{Daily Algebra Drill \#3}{Jane Doe}
% %% AGENTS: When generating a new sheet, ask the human contributor for the name they want credited in argument 2.
% For answer files, pass the "--- Answers \& Investigations" suffix in the arg.
\newcommand{\SpeedTitleBlock}[2]{%
  \begin{center}
    {\LARGE\bfseries #1}\\[4pt]
    {\normalsize\itshape Sheet by #2}\\[6pt]
    \hrule\vspace{4pt}
    \begin{tabular}{llcc}
      \toprule
      \textbf{Section} & \textbf{Topic} & \textbf{Questions} & \textbf{Target time}\\
      \midrule
      A & Rapid Recognition          & 10 & 2 min 30 sec \\
      B & Manipulation Drills        & 10 & 8 min        \\
      C & Substitution \& Structure  & 8  & 10 min       \\
      D & Challenge                  & 5  & 15 min       \\
      \bottomrule
    \end{tabular}
    \vspace{4pt}
    \hrule
  \end{center}
}

% ── Sheet metadata: topic + the day's new tools ────────────────────────────────
% Usage (immediately after \SpeedTitleBlock, sheets only):
%   \SpeedMeta{Expanding, factorising, and the identity toolkit}
%             {difference of squares, sum/product form, completing the square}
%
% Arg 1 = the sheet's topic in one line. The website lists this instead of
%         "Sheet 03", so write the thing a reader would actually search for.
% Arg 2 = comma-separated, at most five: the tools this sheet introduces that
%         earlier sheets in the pillar did not. These become the website's tags.
%         Leave out anything the reader already met — the tags are meant to say
%         what is new today, not to summarise the sheet.
%
% Both are parsed straight out of this file by tools/build_website.py, so the
% sheet stays the single source of truth and the site cannot drift from it.
%
% This renders NOTHING. It is a declaration for the build script, not a piece of
% the sheet's layout: adding it to a sheet must not change that sheet's PDF, so
% the 25 existing sheets can be annotated without a single one of them being
% reissued. If the toolkit should also appear on the page, that is a separate
% decision about the sheet design — make it deliberately, here, not by leaving
% this macro printing something nobody asked it to.
\newcommand{\SpeedMeta}[2]{}

% ── Answer / Method / Investigate-further boxes (answer files only) ────────────
\newmdenv[
  backgroundcolor=lightgray,
  linecolor=medgray,
  linewidth=0.5pt,
  leftmargin=0pt, rightmargin=0pt,
  innerleftmargin=8pt, innerrightmargin=8pt,
  innertopmargin=5pt, innerbottommargin=5pt,
  skipabove=4pt, skipbelow=4pt
]{ansbox}

\newmdenv[
  backgroundcolor=white,
  linecolor=darkgray,
  linewidth=0.8pt,
  leftline=true, rightline=false, topline=false, bottomline=false,
  leftmargin=0pt, rightmargin=0pt,
  innerleftmargin=8pt, innerrightmargin=6pt,
  innertopmargin=4pt, innerbottommargin=4pt,
  skipabove=4pt, skipbelow=6pt
]{invbox}

\newcommand{\ans}[1]{%
  \begin{ansbox}
    \textbf{Answer:} #1
  \end{ansbox}}

\newcommand{\method}[1]{%
  {\small\textbf{Method:} #1}\par}

\newcommand{\inv}[1]{%
  \begin{invbox}
    \textbf{\small Investigate further:} {\small #1}
  \end{invbox}}

% ── Misc helpers ────────────────────────────────────────────────────────────────
\newcommand{\qn}[1]{\textbf{#1.}\;}

% Mistake-linking: attach a Top 5 Patterns / Common Traps bullet to the question
% label(s) in this sheet where it actually shows up, e.g. \seealso{B6, D2}
\newcommand{\seealso}[1]{\hfill\textit{\footnotesize(see #1)}}

% ── Graph options macro ─────────────────────────────────────────────────────────
% Usage: \GraphOptions{tikz code for A}{tikz code for B}{tikz code for C}{tikz code for D}
\newcommand{\GraphOptions}[4]{%
  \begin{center}
    \begin{tabular}{cc}
      \textbf{A)} & \textbf{B)} \\
      #1 & #2 \\[10pt]
      \textbf{C)} & \textbf{D)} \\
      #3 & #4
    \end{tabular}
  \end{center}
}

% ── Closing motivational rule + cross-promo footer (sheets only) ────────────────
% Usage: \SpeedClosing{"Quote text."}
\newcommand{\SpeedClosing}[1]{%
  \vspace{20pt}
  \begin{center}
  \rule{0.6\textwidth}{0.4pt}\\[4pt]
  {\small\itshape #1}\\[4pt]
  \rule{0.6\textwidth}{0.4pt}
  \end{center}
  \begin{center}
  {\small Found a mistake or want to contribute a sheet?}\\
  {\small Visit the open-source project at \href{https://github.com/The-CerealDev/speed-maths}{github.com/The-CerealDev/speed-maths}}
  \end{center}
}
 and before \begin{document}):
%   \SpeedHeader{Algebra}{3}
% First arg  = pillar name shown as "Speed <Pillar> --- Daily Drill"
% Second arg = worksheet number
\pagestyle{fancy}
\newcommand{\SpeedHeader}[2]{%
  \fancyhf{}%
  \renewcommand{\headrulewidth}{0.4pt}%
  \setlength{\headheight}{14pt}%
  \fancyhead[L]{\small\textbf{Speed #1 --- Daily Drill}}%
  \fancyhead[R]{\small Worksheet \##2}%
  \fancyfoot[L]{\small \SpeedExamLine}%
  \fancyfoot[C]{\small\thepage}%
  \fancyfoot[R]{\small \SpeedCredit}%
  \renewcommand{\footrulewidth}{0.4pt}%
}

% ── Title block (used at the top of every sheet AND answer file) ───────────────
% Usage: \SpeedTitleBlock{Daily Algebra Drill \#3}{Jane Doe}
% %% AGENTS: When generating a new sheet, ask the human contributor for the name they want credited in argument 2.
% For answer files, pass the "--- Answers \& Investigations" suffix in the arg.
\newcommand{\SpeedTitleBlock}[2]{%
  \begin{center}
    {\LARGE\bfseries #1}\\[4pt]
    {\normalsize\itshape Sheet by #2}\\[6pt]
    \hrule\vspace{4pt}
    \begin{tabular}{llcc}
      \toprule
      \textbf{Section} & \textbf{Topic} & \textbf{Questions} & \textbf{Target time}\\
      \midrule
      A & Rapid Recognition          & 10 & 2 min 30 sec \\
      B & Manipulation Drills        & 10 & 8 min        \\
      C & Substitution \& Structure  & 8  & 10 min       \\
      D & Challenge                  & 5  & 15 min       \\
      \bottomrule
    \end{tabular}
    \vspace{4pt}
    \hrule
  \end{center}
}

% ── Sheet metadata: topic + the day's new tools ────────────────────────────────
% Usage (immediately after \SpeedTitleBlock, sheets only):
%   \SpeedMeta{Expanding, factorising, and the identity toolkit}
%             {difference of squares, sum/product form, completing the square}
%
% Arg 1 = the sheet's topic in one line. The website lists this instead of
%         "Sheet 03", so write the thing a reader would actually search for.
% Arg 2 = comma-separated, at most five: the tools this sheet introduces that
%         earlier sheets in the pillar did not. These become the website's tags.
%         Leave out anything the reader already met — the tags are meant to say
%         what is new today, not to summarise the sheet.
%
% Both are parsed straight out of this file by tools/build_website.py, so the
% sheet stays the single source of truth and the site cannot drift from it.
%
% This renders NOTHING. It is a declaration for the build script, not a piece of
% the sheet's layout: adding it to a sheet must not change that sheet's PDF, so
% the 25 existing sheets can be annotated without a single one of them being
% reissued. If the toolkit should also appear on the page, that is a separate
% decision about the sheet design — make it deliberately, here, not by leaving
% this macro printing something nobody asked it to.
\newcommand{\SpeedMeta}[2]{}

% ── Answer / Method / Investigate-further boxes (answer files only) ────────────
\newmdenv[
  backgroundcolor=lightgray,
  linecolor=medgray,
  linewidth=0.5pt,
  leftmargin=0pt, rightmargin=0pt,
  innerleftmargin=8pt, innerrightmargin=8pt,
  innertopmargin=5pt, innerbottommargin=5pt,
  skipabove=4pt, skipbelow=4pt
]{ansbox}

\newmdenv[
  backgroundcolor=white,
  linecolor=darkgray,
  linewidth=0.8pt,
  leftline=true, rightline=false, topline=false, bottomline=false,
  leftmargin=0pt, rightmargin=0pt,
  innerleftmargin=8pt, innerrightmargin=6pt,
  innertopmargin=4pt, innerbottommargin=4pt,
  skipabove=4pt, skipbelow=6pt
]{invbox}

\newcommand{\ans}[1]{%
  \begin{ansbox}
    \textbf{Answer:} #1
  \end{ansbox}}

\newcommand{\method}[1]{%
  {\small\textbf{Method:} #1}\par}

\newcommand{\inv}[1]{%
  \begin{invbox}
    \textbf{\small Investigate further:} {\small #1}
  \end{invbox}}

% ── Misc helpers ────────────────────────────────────────────────────────────────
\newcommand{\qn}[1]{\textbf{#1.}\;}

% Mistake-linking: attach a Top 5 Patterns / Common Traps bullet to the question
% label(s) in this sheet where it actually shows up, e.g. \seealso{B6, D2}
\newcommand{\seealso}[1]{\hfill\textit{\footnotesize(see #1)}}

% ── Graph options macro ─────────────────────────────────────────────────────────
% Usage: \GraphOptions{tikz code for A}{tikz code for B}{tikz code for C}{tikz code for D}
\newcommand{\GraphOptions}[4]{%
  \begin{center}
    \begin{tabular}{cc}
      \textbf{A)} & \textbf{B)} \\
      #1 & #2 \\[10pt]
      \textbf{C)} & \textbf{D)} \\
      #3 & #4
    \end{tabular}
  \end{center}
}

% ── Closing motivational rule + cross-promo footer (sheets only) ────────────────
% Usage: \SpeedClosing{"Quote text."}
\newcommand{\SpeedClosing}[1]{%
  \vspace{20pt}
  \begin{center}
  \rule{0.6\textwidth}{0.4pt}\\[4pt]
  {\small\itshape #1}\\[4pt]
  \rule{0.6\textwidth}{0.4pt}
  \end{center}
  \begin{center}
  {\small Found a mistake or want to contribute a sheet?}\\
  {\small Visit the open-source project at \href{https://github.com/The-CerealDev/speed-maths}{github.com/The-CerealDev/speed-maths}}
  \end{center}
}

% Editing THIS file changes the look of every sheet/answer across every pillar.
% ══════════════════════════════════════════════════════════════════════════════

\usepackage[a4paper, top=25mm, bottom=25mm, left=22mm, right=22mm]{geometry}
\usepackage{amsmath, amssymb}
\usepackage{enumitem}
\usepackage{booktabs}
\usepackage{array}
\usepackage{parskip}
\usepackage{microtype}
\usepackage{fancyhdr}
\usepackage{titlesec}
\usepackage{mdframed}
\usepackage{xcolor}
\usepackage[hidelinks]{hyperref}
\usepackage{graphicx}
\usepackage{tikz}
\usepackage{pgfplots}
\pgfplotsset{compat=1.18}

% ── Colours ───────────────────────────────────────────────────────────────────
\definecolor{darkgray}{gray}{0.15}
\definecolor{medgray}{gray}{0.45}
\definecolor{lightgray}{gray}{0.93}
\definecolor{boxgray}{gray}{0.88}

% ── Section formatting ──────────────────────────────────────────────────────────
\titleformat{\section}[block]
  {\normalfont\large\bfseries}{}
  {0pt}{}[\vspace{2pt}\hrule\vspace{6pt}]
\titlespacing*{\section}{0pt}{14pt}{4pt}

% ── List formatting ─────────────────────────────────────────────────────────────
\setlist[enumerate,1]{label=\textbf{\arabic*.}, leftmargin=2em, itemsep=10pt}

% ── Branding constants (edit here, not per-file) ────────────────────────────────
\newcommand{\SpeedExamLine}{\textbf{Speed Maths:} Competition \& Exam Prep}
\newcommand{\SpeedCredit}{Series by \href{https://www.linkedin.com/in/cerealdev/}{David Akinyele-Aje}}

% ── Header / footer ──────────────────────────────────────────────────────────────
% Usage (once, after % main (standalone preamble for editing)
% vim: nowrap: spell:
% ══════════════════════════════════════════════════════════════════════════════
% Speed Maths --- shared preamble
% Included by every sheet and answer file via:  % main (standalone preamble for editing)
% vim: nowrap: spell:
% ══════════════════════════════════════════════════════════════════════════════
% Speed Maths --- shared preamble
% Included by every sheet and answer file via:  \input{../../shared/preamble}
% Editing THIS file changes the look of every sheet/answer across every pillar.
% ══════════════════════════════════════════════════════════════════════════════

\usepackage[a4paper, top=25mm, bottom=25mm, left=22mm, right=22mm]{geometry}
\usepackage{amsmath, amssymb}
\usepackage{enumitem}
\usepackage{booktabs}
\usepackage{array}
\usepackage{parskip}
\usepackage{microtype}
\usepackage{fancyhdr}
\usepackage{titlesec}
\usepackage{mdframed}
\usepackage{xcolor}
\usepackage[hidelinks]{hyperref}
\usepackage{graphicx}
\usepackage{tikz}
\usepackage{pgfplots}
\pgfplotsset{compat=1.18}

% ── Colours ───────────────────────────────────────────────────────────────────
\definecolor{darkgray}{gray}{0.15}
\definecolor{medgray}{gray}{0.45}
\definecolor{lightgray}{gray}{0.93}
\definecolor{boxgray}{gray}{0.88}

% ── Section formatting ──────────────────────────────────────────────────────────
\titleformat{\section}[block]
  {\normalfont\large\bfseries}{}
  {0pt}{}[\vspace{2pt}\hrule\vspace{6pt}]
\titlespacing*{\section}{0pt}{14pt}{4pt}

% ── List formatting ─────────────────────────────────────────────────────────────
\setlist[enumerate,1]{label=\textbf{\arabic*.}, leftmargin=2em, itemsep=10pt}

% ── Branding constants (edit here, not per-file) ────────────────────────────────
\newcommand{\SpeedExamLine}{\textbf{Speed Maths:} Competition \& Exam Prep}
\newcommand{\SpeedCredit}{Series by \href{https://www.linkedin.com/in/cerealdev/}{David Akinyele-Aje}}

% ── Header / footer ──────────────────────────────────────────────────────────────
% Usage (once, after \input{../../shared/preamble} and before \begin{document}):
%   \SpeedHeader{Algebra}{3}
% First arg  = pillar name shown as "Speed <Pillar> --- Daily Drill"
% Second arg = worksheet number
\pagestyle{fancy}
\newcommand{\SpeedHeader}[2]{%
  \fancyhf{}%
  \renewcommand{\headrulewidth}{0.4pt}%
  \setlength{\headheight}{14pt}%
  \fancyhead[L]{\small\textbf{Speed #1 --- Daily Drill}}%
  \fancyhead[R]{\small Worksheet \##2}%
  \fancyfoot[L]{\small \SpeedExamLine}%
  \fancyfoot[C]{\small\thepage}%
  \fancyfoot[R]{\small \SpeedCredit}%
  \renewcommand{\footrulewidth}{0.4pt}%
}

% ── Title block (used at the top of every sheet AND answer file) ───────────────
% Usage: \SpeedTitleBlock{Daily Algebra Drill \#3}{Jane Doe}
% %% AGENTS: When generating a new sheet, ask the human contributor for the name they want credited in argument 2.
% For answer files, pass the "--- Answers \& Investigations" suffix in the arg.
\newcommand{\SpeedTitleBlock}[2]{%
  \begin{center}
    {\LARGE\bfseries #1}\\[4pt]
    {\normalsize\itshape Sheet by #2}\\[6pt]
    \hrule\vspace{4pt}
    \begin{tabular}{llcc}
      \toprule
      \textbf{Section} & \textbf{Topic} & \textbf{Questions} & \textbf{Target time}\\
      \midrule
      A & Rapid Recognition          & 10 & 2 min 30 sec \\
      B & Manipulation Drills        & 10 & 8 min        \\
      C & Substitution \& Structure  & 8  & 10 min       \\
      D & Challenge                  & 5  & 15 min       \\
      \bottomrule
    \end{tabular}
    \vspace{4pt}
    \hrule
  \end{center}
}

% ── Sheet metadata: topic + the day's new tools ────────────────────────────────
% Usage (immediately after \SpeedTitleBlock, sheets only):
%   \SpeedMeta{Expanding, factorising, and the identity toolkit}
%             {difference of squares, sum/product form, completing the square}
%
% Arg 1 = the sheet's topic in one line. The website lists this instead of
%         "Sheet 03", so write the thing a reader would actually search for.
% Arg 2 = comma-separated, at most five: the tools this sheet introduces that
%         earlier sheets in the pillar did not. These become the website's tags.
%         Leave out anything the reader already met — the tags are meant to say
%         what is new today, not to summarise the sheet.
%
% Both are parsed straight out of this file by tools/build_website.py, so the
% sheet stays the single source of truth and the site cannot drift from it.
%
% This renders NOTHING. It is a declaration for the build script, not a piece of
% the sheet's layout: adding it to a sheet must not change that sheet's PDF, so
% the 25 existing sheets can be annotated without a single one of them being
% reissued. If the toolkit should also appear on the page, that is a separate
% decision about the sheet design — make it deliberately, here, not by leaving
% this macro printing something nobody asked it to.
\newcommand{\SpeedMeta}[2]{}

% ── Answer / Method / Investigate-further boxes (answer files only) ────────────
\newmdenv[
  backgroundcolor=lightgray,
  linecolor=medgray,
  linewidth=0.5pt,
  leftmargin=0pt, rightmargin=0pt,
  innerleftmargin=8pt, innerrightmargin=8pt,
  innertopmargin=5pt, innerbottommargin=5pt,
  skipabove=4pt, skipbelow=4pt
]{ansbox}

\newmdenv[
  backgroundcolor=white,
  linecolor=darkgray,
  linewidth=0.8pt,
  leftline=true, rightline=false, topline=false, bottomline=false,
  leftmargin=0pt, rightmargin=0pt,
  innerleftmargin=8pt, innerrightmargin=6pt,
  innertopmargin=4pt, innerbottommargin=4pt,
  skipabove=4pt, skipbelow=6pt
]{invbox}

\newcommand{\ans}[1]{%
  \begin{ansbox}
    \textbf{Answer:} #1
  \end{ansbox}}

\newcommand{\method}[1]{%
  {\small\textbf{Method:} #1}\par}

\newcommand{\inv}[1]{%
  \begin{invbox}
    \textbf{\small Investigate further:} {\small #1}
  \end{invbox}}

% ── Misc helpers ────────────────────────────────────────────────────────────────
\newcommand{\qn}[1]{\textbf{#1.}\;}

% Mistake-linking: attach a Top 5 Patterns / Common Traps bullet to the question
% label(s) in this sheet where it actually shows up, e.g. \seealso{B6, D2}
\newcommand{\seealso}[1]{\hfill\textit{\footnotesize(see #1)}}

% ── Graph options macro ─────────────────────────────────────────────────────────
% Usage: \GraphOptions{tikz code for A}{tikz code for B}{tikz code for C}{tikz code for D}
\newcommand{\GraphOptions}[4]{%
  \begin{center}
    \begin{tabular}{cc}
      \textbf{A)} & \textbf{B)} \\
      #1 & #2 \\[10pt]
      \textbf{C)} & \textbf{D)} \\
      #3 & #4
    \end{tabular}
  \end{center}
}

% ── Closing motivational rule + cross-promo footer (sheets only) ────────────────
% Usage: \SpeedClosing{"Quote text."}
\newcommand{\SpeedClosing}[1]{%
  \vspace{20pt}
  \begin{center}
  \rule{0.6\textwidth}{0.4pt}\\[4pt]
  {\small\itshape #1}\\[4pt]
  \rule{0.6\textwidth}{0.4pt}
  \end{center}
  \begin{center}
  {\small Found a mistake or want to contribute a sheet?}\\
  {\small Visit the open-source project at \href{https://github.com/The-CerealDev/speed-maths}{github.com/The-CerealDev/speed-maths}}
  \end{center}
}

% Editing THIS file changes the look of every sheet/answer across every pillar.
% ══════════════════════════════════════════════════════════════════════════════

\usepackage[a4paper, top=25mm, bottom=25mm, left=22mm, right=22mm]{geometry}
\usepackage{amsmath, amssymb}
\usepackage{enumitem}
\usepackage{booktabs}
\usepackage{array}
\usepackage{parskip}
\usepackage{microtype}
\usepackage{fancyhdr}
\usepackage{titlesec}
\usepackage{mdframed}
\usepackage{xcolor}
\usepackage[hidelinks]{hyperref}
\usepackage{graphicx}
\usepackage{tikz}
\usepackage{pgfplots}
\pgfplotsset{compat=1.18}

% ── Colours ───────────────────────────────────────────────────────────────────
\definecolor{darkgray}{gray}{0.15}
\definecolor{medgray}{gray}{0.45}
\definecolor{lightgray}{gray}{0.93}
\definecolor{boxgray}{gray}{0.88}

% ── Section formatting ──────────────────────────────────────────────────────────
\titleformat{\section}[block]
  {\normalfont\large\bfseries}{}
  {0pt}{}[\vspace{2pt}\hrule\vspace{6pt}]
\titlespacing*{\section}{0pt}{14pt}{4pt}

% ── List formatting ─────────────────────────────────────────────────────────────
\setlist[enumerate,1]{label=\textbf{\arabic*.}, leftmargin=2em, itemsep=10pt}

% ── Branding constants (edit here, not per-file) ────────────────────────────────
\newcommand{\SpeedExamLine}{\textbf{Speed Maths:} Competition \& Exam Prep}
\newcommand{\SpeedCredit}{Series by \href{https://www.linkedin.com/in/cerealdev/}{David Akinyele-Aje}}

% ── Header / footer ──────────────────────────────────────────────────────────────
% Usage (once, after % main (standalone preamble for editing)
% vim: nowrap: spell:
% ══════════════════════════════════════════════════════════════════════════════
% Speed Maths --- shared preamble
% Included by every sheet and answer file via:  \input{../../shared/preamble}
% Editing THIS file changes the look of every sheet/answer across every pillar.
% ══════════════════════════════════════════════════════════════════════════════

\usepackage[a4paper, top=25mm, bottom=25mm, left=22mm, right=22mm]{geometry}
\usepackage{amsmath, amssymb}
\usepackage{enumitem}
\usepackage{booktabs}
\usepackage{array}
\usepackage{parskip}
\usepackage{microtype}
\usepackage{fancyhdr}
\usepackage{titlesec}
\usepackage{mdframed}
\usepackage{xcolor}
\usepackage[hidelinks]{hyperref}
\usepackage{graphicx}
\usepackage{tikz}
\usepackage{pgfplots}
\pgfplotsset{compat=1.18}

% ── Colours ───────────────────────────────────────────────────────────────────
\definecolor{darkgray}{gray}{0.15}
\definecolor{medgray}{gray}{0.45}
\definecolor{lightgray}{gray}{0.93}
\definecolor{boxgray}{gray}{0.88}

% ── Section formatting ──────────────────────────────────────────────────────────
\titleformat{\section}[block]
  {\normalfont\large\bfseries}{}
  {0pt}{}[\vspace{2pt}\hrule\vspace{6pt}]
\titlespacing*{\section}{0pt}{14pt}{4pt}

% ── List formatting ─────────────────────────────────────────────────────────────
\setlist[enumerate,1]{label=\textbf{\arabic*.}, leftmargin=2em, itemsep=10pt}

% ── Branding constants (edit here, not per-file) ────────────────────────────────
\newcommand{\SpeedExamLine}{\textbf{Speed Maths:} Competition \& Exam Prep}
\newcommand{\SpeedCredit}{Series by \href{https://www.linkedin.com/in/cerealdev/}{David Akinyele-Aje}}

% ── Header / footer ──────────────────────────────────────────────────────────────
% Usage (once, after \input{../../shared/preamble} and before \begin{document}):
%   \SpeedHeader{Algebra}{3}
% First arg  = pillar name shown as "Speed <Pillar> --- Daily Drill"
% Second arg = worksheet number
\pagestyle{fancy}
\newcommand{\SpeedHeader}[2]{%
  \fancyhf{}%
  \renewcommand{\headrulewidth}{0.4pt}%
  \setlength{\headheight}{14pt}%
  \fancyhead[L]{\small\textbf{Speed #1 --- Daily Drill}}%
  \fancyhead[R]{\small Worksheet \##2}%
  \fancyfoot[L]{\small \SpeedExamLine}%
  \fancyfoot[C]{\small\thepage}%
  \fancyfoot[R]{\small \SpeedCredit}%
  \renewcommand{\footrulewidth}{0.4pt}%
}

% ── Title block (used at the top of every sheet AND answer file) ───────────────
% Usage: \SpeedTitleBlock{Daily Algebra Drill \#3}{Jane Doe}
% %% AGENTS: When generating a new sheet, ask the human contributor for the name they want credited in argument 2.
% For answer files, pass the "--- Answers \& Investigations" suffix in the arg.
\newcommand{\SpeedTitleBlock}[2]{%
  \begin{center}
    {\LARGE\bfseries #1}\\[4pt]
    {\normalsize\itshape Sheet by #2}\\[6pt]
    \hrule\vspace{4pt}
    \begin{tabular}{llcc}
      \toprule
      \textbf{Section} & \textbf{Topic} & \textbf{Questions} & \textbf{Target time}\\
      \midrule
      A & Rapid Recognition          & 10 & 2 min 30 sec \\
      B & Manipulation Drills        & 10 & 8 min        \\
      C & Substitution \& Structure  & 8  & 10 min       \\
      D & Challenge                  & 5  & 15 min       \\
      \bottomrule
    \end{tabular}
    \vspace{4pt}
    \hrule
  \end{center}
}

% ── Sheet metadata: topic + the day's new tools ────────────────────────────────
% Usage (immediately after \SpeedTitleBlock, sheets only):
%   \SpeedMeta{Expanding, factorising, and the identity toolkit}
%             {difference of squares, sum/product form, completing the square}
%
% Arg 1 = the sheet's topic in one line. The website lists this instead of
%         "Sheet 03", so write the thing a reader would actually search for.
% Arg 2 = comma-separated, at most five: the tools this sheet introduces that
%         earlier sheets in the pillar did not. These become the website's tags.
%         Leave out anything the reader already met — the tags are meant to say
%         what is new today, not to summarise the sheet.
%
% Both are parsed straight out of this file by tools/build_website.py, so the
% sheet stays the single source of truth and the site cannot drift from it.
%
% This renders NOTHING. It is a declaration for the build script, not a piece of
% the sheet's layout: adding it to a sheet must not change that sheet's PDF, so
% the 25 existing sheets can be annotated without a single one of them being
% reissued. If the toolkit should also appear on the page, that is a separate
% decision about the sheet design — make it deliberately, here, not by leaving
% this macro printing something nobody asked it to.
\newcommand{\SpeedMeta}[2]{}

% ── Answer / Method / Investigate-further boxes (answer files only) ────────────
\newmdenv[
  backgroundcolor=lightgray,
  linecolor=medgray,
  linewidth=0.5pt,
  leftmargin=0pt, rightmargin=0pt,
  innerleftmargin=8pt, innerrightmargin=8pt,
  innertopmargin=5pt, innerbottommargin=5pt,
  skipabove=4pt, skipbelow=4pt
]{ansbox}

\newmdenv[
  backgroundcolor=white,
  linecolor=darkgray,
  linewidth=0.8pt,
  leftline=true, rightline=false, topline=false, bottomline=false,
  leftmargin=0pt, rightmargin=0pt,
  innerleftmargin=8pt, innerrightmargin=6pt,
  innertopmargin=4pt, innerbottommargin=4pt,
  skipabove=4pt, skipbelow=6pt
]{invbox}

\newcommand{\ans}[1]{%
  \begin{ansbox}
    \textbf{Answer:} #1
  \end{ansbox}}

\newcommand{\method}[1]{%
  {\small\textbf{Method:} #1}\par}

\newcommand{\inv}[1]{%
  \begin{invbox}
    \textbf{\small Investigate further:} {\small #1}
  \end{invbox}}

% ── Misc helpers ────────────────────────────────────────────────────────────────
\newcommand{\qn}[1]{\textbf{#1.}\;}

% Mistake-linking: attach a Top 5 Patterns / Common Traps bullet to the question
% label(s) in this sheet where it actually shows up, e.g. \seealso{B6, D2}
\newcommand{\seealso}[1]{\hfill\textit{\footnotesize(see #1)}}

% ── Graph options macro ─────────────────────────────────────────────────────────
% Usage: \GraphOptions{tikz code for A}{tikz code for B}{tikz code for C}{tikz code for D}
\newcommand{\GraphOptions}[4]{%
  \begin{center}
    \begin{tabular}{cc}
      \textbf{A)} & \textbf{B)} \\
      #1 & #2 \\[10pt]
      \textbf{C)} & \textbf{D)} \\
      #3 & #4
    \end{tabular}
  \end{center}
}

% ── Closing motivational rule + cross-promo footer (sheets only) ────────────────
% Usage: \SpeedClosing{"Quote text."}
\newcommand{\SpeedClosing}[1]{%
  \vspace{20pt}
  \begin{center}
  \rule{0.6\textwidth}{0.4pt}\\[4pt]
  {\small\itshape #1}\\[4pt]
  \rule{0.6\textwidth}{0.4pt}
  \end{center}
  \begin{center}
  {\small Found a mistake or want to contribute a sheet?}\\
  {\small Visit the open-source project at \href{https://github.com/The-CerealDev/speed-maths}{github.com/The-CerealDev/speed-maths}}
  \end{center}
}
 and before \begin{document}):
%   \SpeedHeader{Algebra}{3}
% First arg  = pillar name shown as "Speed <Pillar> --- Daily Drill"
% Second arg = worksheet number
\pagestyle{fancy}
\newcommand{\SpeedHeader}[2]{%
  \fancyhf{}%
  \renewcommand{\headrulewidth}{0.4pt}%
  \setlength{\headheight}{14pt}%
  \fancyhead[L]{\small\textbf{Speed #1 --- Daily Drill}}%
  \fancyhead[R]{\small Worksheet \##2}%
  \fancyfoot[L]{\small \SpeedExamLine}%
  \fancyfoot[C]{\small\thepage}%
  \fancyfoot[R]{\small \SpeedCredit}%
  \renewcommand{\footrulewidth}{0.4pt}%
}

% ── Title block (used at the top of every sheet AND answer file) ───────────────
% Usage: \SpeedTitleBlock{Daily Algebra Drill \#3}{Jane Doe}
% %% AGENTS: When generating a new sheet, ask the human contributor for the name they want credited in argument 2.
% For answer files, pass the "--- Answers \& Investigations" suffix in the arg.
\newcommand{\SpeedTitleBlock}[2]{%
  \begin{center}
    {\LARGE\bfseries #1}\\[4pt]
    {\normalsize\itshape Sheet by #2}\\[6pt]
    \hrule\vspace{4pt}
    \begin{tabular}{llcc}
      \toprule
      \textbf{Section} & \textbf{Topic} & \textbf{Questions} & \textbf{Target time}\\
      \midrule
      A & Rapid Recognition          & 10 & 2 min 30 sec \\
      B & Manipulation Drills        & 10 & 8 min        \\
      C & Substitution \& Structure  & 8  & 10 min       \\
      D & Challenge                  & 5  & 15 min       \\
      \bottomrule
    \end{tabular}
    \vspace{4pt}
    \hrule
  \end{center}
}

% ── Sheet metadata: topic + the day's new tools ────────────────────────────────
% Usage (immediately after \SpeedTitleBlock, sheets only):
%   \SpeedMeta{Expanding, factorising, and the identity toolkit}
%             {difference of squares, sum/product form, completing the square}
%
% Arg 1 = the sheet's topic in one line. The website lists this instead of
%         "Sheet 03", so write the thing a reader would actually search for.
% Arg 2 = comma-separated, at most five: the tools this sheet introduces that
%         earlier sheets in the pillar did not. These become the website's tags.
%         Leave out anything the reader already met — the tags are meant to say
%         what is new today, not to summarise the sheet.
%
% Both are parsed straight out of this file by tools/build_website.py, so the
% sheet stays the single source of truth and the site cannot drift from it.
%
% This renders NOTHING. It is a declaration for the build script, not a piece of
% the sheet's layout: adding it to a sheet must not change that sheet's PDF, so
% the 25 existing sheets can be annotated without a single one of them being
% reissued. If the toolkit should also appear on the page, that is a separate
% decision about the sheet design — make it deliberately, here, not by leaving
% this macro printing something nobody asked it to.
\newcommand{\SpeedMeta}[2]{}

% ── Answer / Method / Investigate-further boxes (answer files only) ────────────
\newmdenv[
  backgroundcolor=lightgray,
  linecolor=medgray,
  linewidth=0.5pt,
  leftmargin=0pt, rightmargin=0pt,
  innerleftmargin=8pt, innerrightmargin=8pt,
  innertopmargin=5pt, innerbottommargin=5pt,
  skipabove=4pt, skipbelow=4pt
]{ansbox}

\newmdenv[
  backgroundcolor=white,
  linecolor=darkgray,
  linewidth=0.8pt,
  leftline=true, rightline=false, topline=false, bottomline=false,
  leftmargin=0pt, rightmargin=0pt,
  innerleftmargin=8pt, innerrightmargin=6pt,
  innertopmargin=4pt, innerbottommargin=4pt,
  skipabove=4pt, skipbelow=6pt
]{invbox}

\newcommand{\ans}[1]{%
  \begin{ansbox}
    \textbf{Answer:} #1
  \end{ansbox}}

\newcommand{\method}[1]{%
  {\small\textbf{Method:} #1}\par}

\newcommand{\inv}[1]{%
  \begin{invbox}
    \textbf{\small Investigate further:} {\small #1}
  \end{invbox}}

% ── Misc helpers ────────────────────────────────────────────────────────────────
\newcommand{\qn}[1]{\textbf{#1.}\;}

% Mistake-linking: attach a Top 5 Patterns / Common Traps bullet to the question
% label(s) in this sheet where it actually shows up, e.g. \seealso{B6, D2}
\newcommand{\seealso}[1]{\hfill\textit{\footnotesize(see #1)}}

% ── Graph options macro ─────────────────────────────────────────────────────────
% Usage: \GraphOptions{tikz code for A}{tikz code for B}{tikz code for C}{tikz code for D}
\newcommand{\GraphOptions}[4]{%
  \begin{center}
    \begin{tabular}{cc}
      \textbf{A)} & \textbf{B)} \\
      #1 & #2 \\[10pt]
      \textbf{C)} & \textbf{D)} \\
      #3 & #4
    \end{tabular}
  \end{center}
}

% ── Closing motivational rule + cross-promo footer (sheets only) ────────────────
% Usage: \SpeedClosing{"Quote text."}
\newcommand{\SpeedClosing}[1]{%
  \vspace{20pt}
  \begin{center}
  \rule{0.6\textwidth}{0.4pt}\\[4pt]
  {\small\itshape #1}\\[4pt]
  \rule{0.6\textwidth}{0.4pt}
  \end{center}
  \begin{center}
  {\small Found a mistake or want to contribute a sheet?}\\
  {\small Visit the open-source project at \href{https://github.com/The-CerealDev/speed-maths}{github.com/The-CerealDev/speed-maths}}
  \end{center}
}
 and before \begin{document}):
%   \SpeedHeader{Algebra}{3}
% First arg  = pillar name shown as "Speed <Pillar> --- Daily Drill"
% Second arg = worksheet number
\pagestyle{fancy}
\newcommand{\SpeedHeader}[2]{%
  \fancyhf{}%
  \renewcommand{\headrulewidth}{0.4pt}%
  \setlength{\headheight}{14pt}%
  \fancyhead[L]{\small\textbf{Speed #1 --- Daily Drill}}%
  \fancyhead[R]{\small Worksheet \##2}%
  \fancyfoot[L]{\small \SpeedExamLine}%
  \fancyfoot[C]{\small\thepage}%
  \fancyfoot[R]{\small \SpeedCredit}%
  \renewcommand{\footrulewidth}{0.4pt}%
}

% ── Title block (used at the top of every sheet AND answer file) ───────────────
% Usage: \SpeedTitleBlock{Daily Algebra Drill \#3}{Jane Doe}
% %% AGENTS: When generating a new sheet, ask the human contributor for the name they want credited in argument 2.
% For answer files, pass the "--- Answers \& Investigations" suffix in the arg.
\newcommand{\SpeedTitleBlock}[2]{%
  \begin{center}
    {\LARGE\bfseries #1}\\[4pt]
    {\normalsize\itshape Sheet by #2}\\[6pt]
    \hrule\vspace{4pt}
    \begin{tabular}{llcc}
      \toprule
      \textbf{Section} & \textbf{Topic} & \textbf{Questions} & \textbf{Target time}\\
      \midrule
      A & Rapid Recognition          & 10 & 2 min 30 sec \\
      B & Manipulation Drills        & 10 & 8 min        \\
      C & Substitution \& Structure  & 8  & 10 min       \\
      D & Challenge                  & 5  & 15 min       \\
      \bottomrule
    \end{tabular}
    \vspace{4pt}
    \hrule
  \end{center}
}

% ── Sheet metadata: topic + the day's new tools ────────────────────────────────
% Usage (immediately after \SpeedTitleBlock, sheets only):
%   \SpeedMeta{Expanding, factorising, and the identity toolkit}
%             {difference of squares, sum/product form, completing the square}
%
% Arg 1 = the sheet's topic in one line. The website lists this instead of
%         "Sheet 03", so write the thing a reader would actually search for.
% Arg 2 = comma-separated, at most five: the tools this sheet introduces that
%         earlier sheets in the pillar did not. These become the website's tags.
%         Leave out anything the reader already met — the tags are meant to say
%         what is new today, not to summarise the sheet.
%
% Both are parsed straight out of this file by tools/build_website.py, so the
% sheet stays the single source of truth and the site cannot drift from it.
%
% This renders NOTHING. It is a declaration for the build script, not a piece of
% the sheet's layout: adding it to a sheet must not change that sheet's PDF, so
% the 25 existing sheets can be annotated without a single one of them being
% reissued. If the toolkit should also appear on the page, that is a separate
% decision about the sheet design — make it deliberately, here, not by leaving
% this macro printing something nobody asked it to.
\newcommand{\SpeedMeta}[2]{}

% ── Answer / Method / Investigate-further boxes (answer files only) ────────────
\newmdenv[
  backgroundcolor=lightgray,
  linecolor=medgray,
  linewidth=0.5pt,
  leftmargin=0pt, rightmargin=0pt,
  innerleftmargin=8pt, innerrightmargin=8pt,
  innertopmargin=5pt, innerbottommargin=5pt,
  skipabove=4pt, skipbelow=4pt
]{ansbox}

\newmdenv[
  backgroundcolor=white,
  linecolor=darkgray,
  linewidth=0.8pt,
  leftline=true, rightline=false, topline=false, bottomline=false,
  leftmargin=0pt, rightmargin=0pt,
  innerleftmargin=8pt, innerrightmargin=6pt,
  innertopmargin=4pt, innerbottommargin=4pt,
  skipabove=4pt, skipbelow=6pt
]{invbox}

\newcommand{\ans}[1]{%
  \begin{ansbox}
    \textbf{Answer:} #1
  \end{ansbox}}

\newcommand{\method}[1]{%
  {\small\textbf{Method:} #1}\par}

\newcommand{\inv}[1]{%
  \begin{invbox}
    \textbf{\small Investigate further:} {\small #1}
  \end{invbox}}

% ── Misc helpers ────────────────────────────────────────────────────────────────
\newcommand{\qn}[1]{\textbf{#1.}\;}

% Mistake-linking: attach a Top 5 Patterns / Common Traps bullet to the question
% label(s) in this sheet where it actually shows up, e.g. \seealso{B6, D2}
\newcommand{\seealso}[1]{\hfill\textit{\footnotesize(see #1)}}

% ── Graph options macro ─────────────────────────────────────────────────────────
% Usage: \GraphOptions{tikz code for A}{tikz code for B}{tikz code for C}{tikz code for D}
\newcommand{\GraphOptions}[4]{%
  \begin{center}
    \begin{tabular}{cc}
      \textbf{A)} & \textbf{B)} \\
      #1 & #2 \\[10pt]
      \textbf{C)} & \textbf{D)} \\
      #3 & #4
    \end{tabular}
  \end{center}
}

% ── Closing motivational rule + cross-promo footer (sheets only) ────────────────
% Usage: \SpeedClosing{"Quote text."}
\newcommand{\SpeedClosing}[1]{%
  \vspace{20pt}
  \begin{center}
  \rule{0.6\textwidth}{0.4pt}\\[4pt]
  {\small\itshape #1}\\[4pt]
  \rule{0.6\textwidth}{0.4pt}
  \end{center}
  \begin{center}
  {\small Found a mistake or want to contribute a sheet?}\\
  {\small Visit the open-source project at \href{https://github.com/The-CerealDev/speed-maths}{github.com/The-CerealDev/speed-maths}}
  \end{center}
}

\SpeedHeader{Combinatorics}{2}

\begin{document}

\SpeedTitleBlock{Daily Combinatorics Drill \#2 --- Answers \& Investigations}
\noindent\textit{New toolkit today: Combinations $\binom{n}{r}$ $\cdot$ Choosing vs Arranging $\cdot$ Complementary Counting $\cdot$ Committee Constructions.}

\section*{Section A \quad Rapid Recognition \hfill{\normalfont\small($\leq$15\,s each)}}
%═══════════════════════════════════════════════════════════════════════════════

\begin{enumerate}

%── A1 ──────────────────────────────────────────────────────────────────────────
\item Evaluate: $\;\dbinom{7}{2}$.

\ans{$21$}
\method{$\dbinom{7}{2}=\dfrac{7\cdot6}{2}=21$. For $\binom{n}{2}$, always: half of $n(n-1)$.}
\inv{$\binom{n}{2}$ is the $(n-1)$th triangular number. Why must the product of two consecutive integers always be even, making $\binom{n}{2}$ an integer? This integrality question generalises to all $\binom{n}{r}$ --- see B-section of sheet 6.}

%── A2 ──────────────────────────────────────────────────────────────────────────
\item Evaluate: $\;\dbinom{8}{3}$.

\ans{$56$}
\method{$\dbinom{8}{3}=\dfrac{8\cdot7\cdot6}{3!}=\dfrac{336}{6}=56$. Cancel before multiplying: $\frac{8\cdot7\cdot6}{6}=8\cdot7$.}
\inv{Note $\binom{8}{3}=56=\binom{8}{5}$. Prove the symmetry $\binom{n}{r}=\binom{n}{n-r}$ \emph{without} algebra: choosing which 3 objects to take is the same decision as choosing which 5 to leave behind.}

%── A3 ──────────────────────────────────────────────────────────────────────────
\item Evaluate: $\;\dbinom{9}{7}$.

\ans{$36$}
\method{Use symmetry first: $\dbinom{9}{7}=\dbinom{9}{2}=\dfrac{9\cdot8}{2}=36$. Never compute a large lower index directly.}
\inv{How large can the saving get? Compare the work in computing $\binom{100}{98}$ directly against $\binom{100}{2}$. Under exam time pressure, the symmetry flip is the single most-used binomial reflex.}

%── A4 ──────────────────────────────────────────────────────────────────────────
\item How many committees of 3 can be chosen from 6 people?

\ans{$\dbinom{6}{3}=20$}
\method{A committee is unordered: $\dbinom{6}{3}=\dfrac{6\cdot5\cdot4}{6}=20$.}
\inv{Contrast with sheet 1's medal question: 3 \emph{distinct} medals from 8 runners was $8\cdot7\cdot6$. Divide-by-$r!$ is exactly the bridge between the two: $\binom{n}{r}=\frac{{}^nP_r}{r!}$. When does a question secretly order the chosen objects?}

%── A5 ──────────────────────────────────────────────────────────────────────────
\item Solve for $n$: $\;\dbinom{n}{2}=45$.

\ans{$n=10$}
\method{$\dfrac{n(n-1)}{2}=45\Rightarrow n(n-1)=90=10\cdot9\Rightarrow n=10$. Spot consecutive-integer products instantly.}
\inv{Which totals $T$ make $\binom{n}{2}=T$ solvable? Exactly the triangular numbers $1,3,6,10,15,\ldots$ Investigate: can a number be triangular in two different ways?}

%── A6 ──────────────────────────────────────────────────────────────────────────
\item How many subsets of a 5-element set have at least 4 elements?

\ans{$\dbinom{5}{4}+\dbinom{5}{5}=5+1=6$}
\method{``At least 4'' means sizes 4 or 5. Two small cases --- add them.}
\inv{``At least'' with \emph{few} qualifying cases: add directly. ``At least'' with \emph{many} qualifying cases: subtract the complement. Recount this both ways ($32-\binom{5}{0}-\binom{5}{1}-\binom{5}{2}-\binom{5}{3}=6$) and note which was faster --- then decide the crossover rule for yourself.}

%── A7 ──────────────────────────────────────────────────────────────────────────
\item Evaluate: $\;\dbinom{6}{0}+\dbinom{6}{1}+\cdots+\dbinom{6}{6}$.

\ans{$2^6=64$}
\method{Row sum of Pascal's triangle: every subset of a 6-set has \emph{some} size, so the sizes partition the $2^6$ subsets.}
\inv{This double count (subsets by size vs.\ in/out choices) is your first serious ``count the same thing two ways'' argument --- the headline tool of sheet 6. Also evaluate the alternating version $\binom{6}{0}-\binom{6}{1}+\binom{6}{2}-\cdots$ (answer 0) and explain it by pairing subsets.}

%── A8 ──────────────────────────────────────────────────────────────────────────
\item How many ways can 2 letters be chosen from $\{$A, B, C, D, E$\}$?

\ans{$\dbinom{5}{2}=10$}
\method{Unordered pairs from 5: $\frac{5\cdot4}{2}=10$.}
\inv{List all 10 pairs quickly in lexicographic order (AB, AC, \ldots, DE). Systematic listing is worth practising: it is the verification tool of last resort when a formula feels doubtful in an exam.}

%── A9 ──────────────────────────────────────────────────────────────────────────
\item True or false: $\;\dbinom{100}{98}=4950$.

\ans{True}
\method{$\dbinom{100}{98}=\dbinom{100}{2}=\dfrac{100\cdot99}{2}=4950$.}
\inv{TMUA-style statement checks reward one clean disproof or one clean identity. Decide true/false: $\binom{n}{2}+\binom{n+1}{2}=n^2$. (True --- prove it, then interpret it as two staircase halves of a square.)}

%── A10 ─────────────────────────────────────────────────────────────────────────
\item A lottery ticket selects 4 numbers from $1$--$10$ (order irrelevant). How many different tickets are there?

\ans{$\dbinom{10}{4}=210$}
\method{Pure choice, no order: $\dfrac{10\cdot9\cdot8\cdot7}{24}=210$.}
\inv{The UK Lotto uses 6 numbers from 59: $\binom{59}{6}=45\,057\,474$. Estimate it mentally as $\frac{59^6}{720}$ to one significant figure and compare --- order-of-magnitude fluency beats exact arithmetic here.}

\end{enumerate}

%═══════════════════════════════════════════════════════════════════════════════
\section*{Section B \quad Manipulation Drills \hfill{\normalfont\small($\leq$50\,s each)}}
%═══════════════════════════════════════════════════════════════════════════════

\begin{enumerate}

%── B1 ──────────────────────────────────────────────────────────────────────────
\item Solve for $n$: $\;\dbinom{n}{3}=35$.

\ans{$n=7$}
\method{$n(n-1)(n-2)=35\cdot6=210=7\cdot6\cdot5$. Three consecutive integers --- read off $n=7$.}
\inv{Estimate first: $n(n-1)(n-2)\approx n^3$, so $n\approx\sqrt[3]{210}\approx6$; test 6 and 7. The estimate-then-check loop solves every $\binom{n}{r}=T$ equation faster than algebra.}

%── B2 ──────────────────────────────────────────────────────────────────────────
\item From 8 men and 6 women, how many committees of 2 men and 2 women can be formed?

\ans{$\dbinom{8}{2}\dbinom{6}{2}=28\cdot15=420$}
\method{Independent choices multiply: choose the men, then the women.}
\inv{Why is it \emph{not} $\binom{14}{4}$ restricted afterwards? Compute how many of the $\binom{14}{4}=1001$ unrestricted committees have the 2--2 split, and check you recover 420. Structure-first beats filter-later.}

%── B3 ──────────────────────────────────────────────────────────────────────────
\item Twelve people at a party all shake hands with each other exactly once. How many handshakes take place?

\ans{$\dbinom{12}{2}=66$}
\method{A handshake \emph{is} an unordered pair of people.}
\inv{Recount by degrees: each of 12 people shakes 11 hands, and $12\cdot11$ counts every handshake twice --- so $66$. This degree-sum argument is the Handshake Lemma; it returns in sheet 6 and in graph theory proper.}

%── B4 ──────────────────────────────────────────────────────────────────────────
\item Ten points, no three collinear: how many segments, and how many triangles?

\ans{$\dbinom{10}{2}=45$ segments;\quad $\dbinom{10}{3}=120$ triangles}
\method{A segment is a pair of points; a triangle is a triple. The no-three-collinear condition guarantees every triple genuinely forms a triangle.}
\inv{Now let exactly 4 of the 10 points be collinear. How many triangles survive? ($120-\binom{4}{3}=116$.) Degenerate cases are subtracted, not case-split --- a complementary-counting reflex.}

%── B5 ──────────────────────────────────────────────────────────────────────────
\item Team of 5 from 11, then a captain from the team. Count both ways.

\ans{$\dbinom{11}{5}\cdot5=2310=11\cdot\dbinom{10}{4}$}
\method{Team first: $462\cdot5$. Captain first: 11 choices, then 4 teammates from the remaining 10: $11\cdot210$. Both give 2310.}
\inv{You have just proved the identity $r\binom{n}{r}=n\binom{n-1}{r-1}$ by counting one committee-with-chair two ways. Write the algebraic proof too, and keep both --- ``committee--chair'' arguments prove many binomial identities without algebra.}

%── B6 ──────────────────────────────────────────────────────────────────────────
\item How many binary strings of length 8 contain exactly three 1s?

\ans{$\dbinom{8}{3}=56$}
\method{A string is determined by \emph{which} 3 of the 8 positions hold the 1s.}
\inv{``Choose the positions'' converts string problems into subset problems. How many length-8 strings have more 1s than 0s? (Sum $\binom{8}{5}+\binom{8}{6}+\binom{8}{7}+\binom{8}{8}=93$.) Could symmetry have halved the work? Careful --- the middle term $\binom{8}{4}$ belongs to neither side.}

%── B7 ──────────────────────────────────────────────────────────────────────────
\item From the 5 vowels and 21 consonants, how many 3-letter \emph{sets} have exactly one vowel and two consonants?

\ans{$\dbinom{5}{1}\dbinom{21}{2}=5\cdot210=1\,050$}
\method{Choose the vowel, then the pair of consonants --- independent choices multiply.}
\inv{Now count 3-letter \emph{strings} (ordered, no repeats) with exactly one vowel: multiply by the $3$ positions for the vowel and $2!$ orders of the consonants --- or by $3!$ on the whole set. Ordered $=$ unordered $\times$ arrangements, always.}

%── B8 ──────────────────────────────────────────────────────────────────────────
\item Solve for $n$: $\;\dbinom{n}{2}+\dbinom{n}{1}=21$.

\ans{$n=6$}
\method{$\frac{n(n-1)}{2}+n=\frac{n(n+1)}{2}=21\Rightarrow n(n+1)=42=6\cdot7\Rightarrow n=6$.}
\inv{You used Pascal's rule backwards: $\binom{n}{2}+\binom{n}{1}=\binom{n+1}{2}$. State Pascal's rule $\binom{n}{r}+\binom{n}{r+1}=\binom{n+1}{r+1}$ and prove it by asking whether a distinguished element is in the chosen set --- sheet 4 builds the whole triangle from this.}

%── B9 ──────────────────────────────────────────────────────────────────────────
\item Committees of 4 from 9 people: how many include Priya? Exclude her?

\ans{Include: $\dbinom{8}{3}=56$;\quad exclude: $\dbinom{8}{4}=70$}
\method{Include her: she is in, choose 3 companions from 8. Exclude her: choose all 4 from the other 8.}
\inv{$56+70=126=\binom{9}{4}$ --- every committee either contains Priya or doesn't. This in/out split \emph{is} Pascal's rule with faces on it. Use the same split to compute $\binom{10}{5}$ from row 9.}

%── B10 ─────────────────────────────────────────────────────────────────────────
\item Routes from $(0,0)$ to $(4,3)$ stepping only right or up?

\ans{$\dbinom{7}{3}=35$}
\method{Every route is a string of 4 R's and 3 U's: choose which 3 of the 7 steps are U.}
\inv{Lattice paths are strings are subsets --- one idea, three costumes. How many routes pass through the point $(2,1)$? (Multiply the two stages: $\binom{3}{1}\binom{4}{2}=18$.) How many \emph{avoid} it?}

\end{enumerate}

%═══════════════════════════════════════════════════════════════════════════════
\section*{Section C \quad Substitution \& Structure \hfill{\normalfont\small($\leq$90\,s each)}}
%═══════════════════════════════════════════════════════════════════════════════

\begin{enumerate}

%── C1 ──────────────────────────────────────────────────────────────────────────
\item Committee of 5 from 7 women and 6 men, at least 4 women.

\ans{$\dbinom{7}{4}\dbinom{6}{1}+\dbinom{7}{5}=35\cdot6+21=231$}
\method{Exactly 4 women or exactly 5 women --- two disjoint cases, then add.}
\inv{A classic overcounting trap lurks here: ``choose 4 women to guarantee the quota ($\binom{7}{4}$), then anyone from the remaining 9'' gives $35\cdot9=315\neq231$. Find the committees counted twice by that argument --- diagnosing it inoculates you permanently.}

%── C2 ──────────────────────────────────────────────────────────────────────────
\item How many four-digit numbers have at least one repeated digit?

\ans{$9000-9\cdot9\cdot8\cdot7=9000-4536=4464$}
\method{Complement: all-distinct four-digit numbers number $9\cdot9\cdot8\cdot7$ (sheet 1 technique). Subtract from 9000.}
\inv{What fraction is that? Nearly half. Estimate the same fraction for 6-digit numbers, then investigate the birthday-paradox structure: with $k$ digits drawn from 10, repeats dominate surprisingly early.}

%── C3 ──────────────────────────────────────────────────────────────────────────
\item How many diagonals does a regular 12-gon have?

\ans{$\dbinom{12}{2}-12=66-12=54$}
\method{Every pair of vertices gives a segment; remove the 12 sides. (Regularity is irrelevant --- only convexity matters.)}
\inv{Derive the general formula $\frac{n(n-3)}{2}$ two ways: (i) $\binom{n}{2}-n$, (ii) each vertex sends $n-3$ diagonals and each is counted twice. Two derivations agreeing is your built-in error check.}

%── C4 ──────────────────────────────────────────────────────────────────────────
\item Triangles from 8 points on a circle; how many avoid point $A$?

\ans{$\dbinom{8}{3}=56$ in total;\quad $\dbinom{7}{3}=35$ avoid $A$}
\method{All triples work (no three points on a circle are collinear). Avoiding $A$: all three vertices come from the other 7 points. Alternatively by complement: $56-\binom{7}{2}=56-21=35$ --- subtract the triangles \emph{through} $A$.}
\inv{Check by symmetry: each of the 56 triangles has 3 vertices, so vertex-uses total $56\cdot3=168=8\cdot21$ --- each point is used 21 times, avoided 35 times, and $21+35=56$. Counting incidences two ways strikes again.}

%── C5 ──────────────────────────────────────────────────────────────────────────
\item Teams of 6 from 6 boys and 6 girls with unequal numbers of each.

\ans{$\dbinom{12}{6}-\dbinom{6}{3}^2=924-400=524$}
\method{Complement: the only \emph{equal} split for a team of 6 is 3--3, counted by $\binom{6}{3}\binom{6}{3}=400$. Subtract from all $\binom{12}{6}=924$ teams. Listing the unequal splits (6--0, 5--1, 4--2 and mirrors) works too but is six cases instead of one subtraction.}
\inv{The identity behind the complement: $\sum_{k}\binom{6}{k}\binom{6}{6-k}=\binom{12}{6}$ --- Vandermonde, proved by asking how many boys ended up in an unrestricted team. Verify the 6-case direct sum gives 524 and feel the difference in effort.}

%── C6 ──────────────────────────────────────────────────────────────────────────
\item 5-card hands from 52 cards containing at least three aces.

\ans{$\dbinom{4}{3}\dbinom{48}{2}+\dbinom{48}{1}=4512+48=4\,560$}
\method{Two disjoint cases. \emph{Exactly three aces:} choose which 3 of the 4 aces, then 2 non-aces: $4\cdot1128$. \emph{All four aces:} the aces are forced, only the fifth card is free: $48$.}
\inv{Why $\binom{4}{3}$ in one case but no $\binom{4}{4}$ anxiety in the other? Forced choices contribute a factor of 1. Extend: compute ``at least two aces'' and check all five ace-counts (0 through 4) sum to $\binom{52}{5}$.}

%── C7 ──────────────────────────────────────────────────────────────────────────
\item A convex $n$-gon has exactly 27 diagonals. Find $n$.

\ans{$n=9$}
\method{$\frac{n(n-3)}{2}=27\Rightarrow n(n-3)=54=9\cdot6\Rightarrow n=9$. (Product of two integers differing by 3 --- read it off.)}
\inv{Which diagonal counts are achievable? Compute $\frac{n(n-3)}{2}$ for $n=4,\ldots,10$: $2,5,9,14,20,27,35$. Investigate: can two different polygons ever have the same diagonal count? (The sequence is strictly increasing --- why does that settle it?)}

%── C8 ──────────────────────────────────────────────────────────────────────────
\item Subsets of $\{1,\ldots,10\}$ containing at least one of 1 and 2.

\ans{$2^{10}-2^8=1024-256=768$}
\method{Complement: subsets avoiding \emph{both} 1 and 2 are exactly the subsets of the remaining 8 elements: $2^8$.}
\inv{Recount by inclusion--exclusion: (contain 1) $+$ (contain 2) $-$ (contain both) $=512+512-256=768$. You have met inclusion--exclusion informally --- sheet 5 makes it a machine. Which route generalises better to ``at least one of 1, 2, 3''?}

\end{enumerate}

%═══════════════════════════════════════════════════════════════════════════════
\section*{Section D \quad Challenge \hfill{\normalfont\small(SMC difficulty)}}
%═══════════════════════════════════════════════════════════════════════════════

\begin{enumerate}

%── D1 ──────────────────────────────────────────────────────────────────────────
\item Rectangles on a $4\times4$ grid. \textit{\small(extends Sheet 1, D2)}

\ans{$\dbinom{5}{2}^2=100$}
\method{A rectangle is determined by choosing 2 of the 5 vertical gridlines and 2 of the 5 horizontal gridlines: $\binom{5}{2}\cdot\binom{5}{2}=10\cdot10$. Sheet 1 counted the spans by listing $(4+3+2+1)$; the binomial coefficient \emph{is} that sum.}
\inv{So $1+2+3+4=\binom{5}{2}$ --- triangular numbers are $\binom{n+1}{2}$ in disguise. Count the \emph{squares} on the same grid ($16+9+4+1=30$), then find how many rectangles on an $8\times8$ board are NOT squares.}

%── D2 ──────────────────────────────────────────────────────────────────────────
\item Subsets of $\{1,\ldots,8\}$ with no two consecutive integers (empty set included). \textit{\small(cf.\ MAT 2023 Q5)}

\ans{$55$}
\method{Let $a_n$ count such subsets of $\{1,\ldots,n\}$. Split on whether $n$ is used: if not, $a_{n-1}$ subsets; if yes, $n-1$ is banned, leaving $a_{n-2}$. So $a_n=a_{n-1}+a_{n-2}$ with $a_1=2$, $a_2=3$ --- Fibonacci! The chain runs $2,3,5,8,13,21,34,55$, so $a_8=55$.}
\inv{Alternative proof: a size-$k$ subset with no neighbours corresponds to choosing $k$ from $8-k+1$ positions after shrinking the gaps, giving $\sum_k\binom{9-k}{k}=55$ --- a Fibonacci number as a diagonal sum in Pascal's triangle. Both proofs preview sheet 7's recurrence toolkit.}

%── D3 ──────────────────────────────────────────────────────────────────────────
\item Committees of any size from 6 juniors and 6 seniors with strictly more juniors.

\ans{$\dfrac{2^{12}-\binom{12}{6}}{2}=\dfrac{4096-924}{2}=1\,586$}
\method{Every committee has more juniors, more seniors, or a tie. Pair each junior with a senior once and for all; swapping every member for their partner is a bijection between ``more juniors'' committees and ``more seniors'' committees, so those two counts are equal. Ties (equal numbers, $k$ of each) number $\sum_k\binom{6}{k}^2=\binom{12}{6}=924$ by Vandermonde. Hence each majority count is $\frac{4096-924}{2}=1586$. No case-by-case summation needed.}
\inv{Case-sum check: $\sum_{j>s}\binom{6}{j}\binom{6}{s}$ really is 1586 --- but the bijection argument scales to 100 juniors and 100 seniors untouched, while the sum does not. Warning: with \emph{unequal} club sizes (7 juniors, 5 seniors) the swap bijection dies. What replaces it? (Nothing clean --- that asymmetric version is a genuine slog, which is exactly why the symmetric version is the beautiful one.)}

%── D4 ──────────────────────────────────────────────────────────────────────────
\item Six people each shake hands with exactly two others. In how many ways?

\ans{$70$}
\method{Everyone having exactly two handshakes forces the handshake pattern to be a union of cycles of length $\geq3$ covering all 6 people: either one 6-cycle or two 3-cycles. \emph{One 6-cycle:} arrange 6 people in a cycle: $\frac{5!}{2}=60$ (rotations and reflection give the same handshake set). \emph{Two 3-cycles:} split into two trios: $\frac{1}{2}\binom{6}{3}=10$, each trio shaking hands in exactly one way. Total $60+10=70$.}
\inv{Why is a ``cycle'' forced? Prove: a graph where every vertex has degree 2 is a disjoint union of cycles. Then redo the count for 7 people (answer: $\frac{6!}{2}+\binom{7}{3}\cdot 1\cdot\frac{3!}{2}=360+35\cdot3=465$ --- careful: a 3-cycle plus a 4-cycle. Work it out and brute-check it).}

%── D5 ──────────────────────────────────────────────────────────────────────────
\item Ten points on a circle, all chords drawn, no three concurrent inside. Interior intersection points?

\ans{$\dbinom{10}{4}=210$}
\method{The key insight: every interior intersection comes from exactly one set of 4 points on the circle (the two crossing chords are the diagonals of the quadrilateral those 4 points form), and every 4 points create exactly one interior crossing. So intersections $\leftrightarrow$ 4-subsets: $\binom{10}{4}=210$. No coordinate geometry, no counting chords.}
\inv{This is a \emph{bijection} doing all the work --- sheet 7's crown tool. Follow-up: how many \emph{triangles} are formed whose vertices are all interior intersection points... is much harder; instead count the regions the chords cut the disc into for $n=1,\ldots,6$ points ($1,2,4,8,16,31$!) and investigate why the doubling pattern breaks.}

\end{enumerate}

%═══════════════════════════════════════════════════════════════════════════════
\section*{Top 5 Patterns Today}
%═══════════════════════════════════════════════════════════════════════════════

\begin{enumerate}[leftmargin=2em, itemsep=4pt]
  \item \textbf{Choose vs arrange.} $\binom{n}{r}=\frac{{}^nP_r}{r!}$: divide by $r!$ exactly when the selection is unordered. Ask ``does swapping two chosen objects give a new outcome?''
  \item \textbf{Symmetry flip.} $\binom{n}{r}=\binom{n}{n-r}$: always compute the smaller index. $\binom{100}{98}$ is a two-second question.
  \item \textbf{Complementary counting.} ``At least one X'' $=$ total $-$ ``no X''. The complement of a hard condition is usually one clean product.
  \item \textbf{Committee--chair double counting.} $r\binom{n}{r}=n\binom{n-1}{r-1}$: proving identities by counting a scenario two ways beats algebra.
  \item \textbf{Objects $\leftrightarrow$ subsets.} Binary strings, lattice paths, handshakes, grid rectangles, chord crossings --- all are ``choose a subset'' in costume. Find the subset.
\end{enumerate}

%═══════════════════════════════════════════════════════════════════════════════
\section*{Common Traps to Avoid}
%═══════════════════════════════════════════════════════════════════════════════

\begin{enumerate}[leftmargin=2em, itemsep=4pt]
  \item \textbf{Ordering an unordered selection.} A committee counted with ${}^nP_r$ is overcounted by a factor of $r!$. Decide ordered/unordered \emph{before} touching numbers.
  \item \textbf{The ``guarantee the quota'' overcount.} ``Choose 4 women, then anyone else'' double-counts committees with 5 women (C1). Quota problems need exact cases or a complement --- never ``reserve then top up''.
  \item \textbf{Adding overlapping cases.} Cases must be \emph{disjoint} before you add. ``Contains 1'' and ``contains 2'' overlap; either subtract the overlap or complement.
  \item \textbf{Forgetting forced choices are free.} All four aces in the hand is a factor of 1, not $\binom{4}{4}$ worth of anxiety --- but \emph{exactly} three aces needs $\binom{4}{3}$.
  \item \textbf{Dividing by symmetry that isn't there.} Handshake circles divide by rotations \emph{and} reflections ($\frac{5!}{2}$); labelled pair-splits divide by $2$ only when the two groups are indistinguishable. Name the symmetry before dividing by it.
\end{enumerate}

\SpeedClosing{``When a count splits into cases, count the complement first ---
the opposite of `at least one' is always simpler.''}

\end{document}
